\chapter{理论与工程约化：从一般动力学到离散网络}
\label{sec:section_5204}
% Compatibility anchors: old section references resolve to this chapter.
\label{sec:tpug_operator_design_formulas}

我们通过离散化将一般状态动力学连接到可执行网络。离散模型必须说明所保留的性质与近似误差，而不是仅凭名称宣称与连续模型精确等价。
\section{图模型的构造}
设 $W=W^\top\ge0$，$L=D-W$。对节点活动 $x$，定义
\begin{equation}\dot x=-Lx-\nabla V(x)+Bu.\end{equation}
固定输入为零时，该式是目标函数 $J(x)=x^\top Lx/2+V(x)$ 的梯度流，所以 $\dot J=-\|\nabla J\|^2\le0$。有输入时需加入 $\nabla J^\top Bu$，不能继续无条件宣称下降。
\section{连续近似的条件}
若图来自空间采样，边权按指定规则近似微分算子，可研究采样尺度趋零时的算子一致性。边界、权重尺度和采样分布都会影响极限。任意知识图不自动对应某个连续流形上的拉普拉斯。
\section{实现与复杂度}
稀疏图更新的单步成本与边数和特征维数相关；稠密注意力仍需要序列长度的二次交互。并行或模拟实现可以改变延迟与常数，但不能免除硬件数量、通信、精度和输入输出成本。
\section{工程不变量与实验}
我们先验证维度、边界、掩码与数值稳定，再检查状态范数、目标下降和任务表现。每项保证只在对应条件内成立。参数归一化可提供范围约束，谱约束可限制局部增益，二者均不能独立保证整个智能体正确行动。
